\section{System Overview and Design}
\label{sec:design}
Application predicts Slack, the Router owns request execution, and workers
independently copy sealed KV to Peer-HBM and Store DRAM (Figure~\ref{fig:system}).
Algorithms 2/3 specify revision 3; the deployed prototype and workload results
use the earlier revision (Section~\ref{sec:implementation}).

\begin{figure*}[tp]
\centering\includegraphics[width=\textwidth]{figures/system-overview.pdf}
\caption{Revision-3 architecture: independent source-owned HBM/DRAM paths
(blue/orange), with metadata in purple. Each tier reserves and acknowledges
separately. Earlier native policy resides in Store Master; revision 3 has an
event reference. Backup, staging and recovery-primary memory are separately
charged. SSD is excluded.}
\Description{Application Slack and Router ownership feed a protection controller;
source workers independently copy to HBM and DRAM, subject to allocator and
shared-path resource limits.}
\label{fig:system}
\end{figure*}

\subsection{Failure model and compatibility}
Fail-stop worker KV loss leaves Application, Router and Controller alive.
Replicas must match model, KV ABI/dtype and parallel layout, with every shard.
Rendezvous hashing binds compatible anti-colocated primary/standby workers;
source failure fences the route before recovery. Evaluated configurations use
TP=PP=1; partitions, control-plane failure and cross-layout recovery are outside
scope. Backup HBM must be copied into separate recovery-primary allocations.

\subsection{Algorithm 1: causal residual-workflow Slack}
The Application updates the causally visible residual DAG at request/tool
returns. GBDT predicts remaining active-node and full known-unstarted-node
durations; structural nodes have zero weight. Forward/backward CPM yields
$s_i=\pos{LS_i-ES_i-m}$ for calibrated margin $m$. Tickets age as
$s_i(t)=\pos{s_i-(t-t_i)}$. Unknown branches are excluded; missing structure or
predictions invalidate active sibling tickets to zero. Calibration uses
disjoint workflows. Predicted float does not guarantee future Slack; graph
processing is $O(|V|+|E|)$ with a fixed predictor.
% Source: Mooncake application/algorithm1.py:update_slack, inspected 2026-09-06.
\begin{algorithm}[t]
\caption{Causal GBDT--CPM update with calibrated Slack}
\label{alg:slack-update}
\begin{tabular}{@{}r@{\hspace{0.5em}}>{\raggedright\arraybackslash}p{\dimexpr\columnwidth-2.5em\relax}@{}}
\multicolumn{2}{@{}>{\raggedright\arraybackslash}p{\columnwidth}@{}}{\textit{Input:} event $e$, time $t$, causal predictor $f_\theta$, workflow-disjoint margin $m\ge0$. $A$: active nodes; $v$: event version.}\\
\multicolumn{2}{@{}>{\raggedright\arraybackslash}p{\columnwidth}@{}}{\textbf{UpdateSlack}$(e,t)$}\\
1 & $(G=(V,E),A,\phi,v)\gets\operatorname{Observe}(e,t)$; $\pi\gets\operatorname{Topo}(G)$.\\
2 & \textbf{if} graph or node metadata is incomplete: publish invalid zero tickets for all $i\in A$; \textbf{return}.\\
3 & For each $u\in V$, set $w_u\gets f_\theta(\phi_u)$ for agent nodes and $w_u\gets0$ for control nodes.\\
4 & \textbf{if} any prediction is unavailable: publish invalid zero tickets for all $i\in A$; \textbf{return}.\\
5 & For $u\in\pi$, set $ES_u\gets0$ at roots, otherwise $ES_u\gets\max_{z\in\operatorname{pred}(u)}(ES_z+w_z)$.\\
6 & $C\gets\max_{u:\operatorname{succ}(u)=\varnothing}(ES_u+w_u)$.\\
7 & For $u\in\operatorname{rev}(\pi)$, set $LS_u\gets C-w_u$ at sinks, otherwise $LS_u\gets\min_{z\in\operatorname{succ}(u)}LS_z-w_u$.\\
8 & For each $i\in A$, set $s_i^{raw}\gets[LS_i-ES_i]^+$ and $s_i\gets[s_i^{raw}-m]^+$.\\
9 & Publish valid tickets $(i,s_i,s_i^{raw},t,v)$ with prediction and calibration provenance.\\
\end{tabular}
\end{algorithm}


\subsection{Algorithm 2: current prefix value and group admission}
\paragraph{Credit is not ordered by insertion.}
The earlier key $H_c=L+A_c/b_c$ can decrease across insertions: $(L,A/b)=(0,10)$
and $(1,1)$ give 10 and 2. Revision 3 removes historical credit and evaluates
current recovery loss under prefix dependencies.

\paragraph{Objects, consumers and value.}
Immutable completed page groups bind parent hash, tokens, tenant, model,
adapter, ABI/layout and non-token identity~\cite{vllmPrefixCaching,zheng2024sglang}.
Committed set $C_r$ is parent-closed; $f_i(C_r)$ is active call $i$'s compatible
frontier. Shared bytes count once; each unique call contributes once to
\begin{equation}
\Phi_r(C,t)=\sum_{i\in\mathcal A(t)}w_i
 \pos{T_i^r(f_i(C),t)-s_i(t)}.
\end{equation}
$T_i^r$ minimizes loading/recompute/replay cost including full recomputation;
$w_i=1$ by default. Ended rounds retain only evictable zero-weight TTL references;
future calls receive no value. This tier-local proxy is not DAG makespan or
long-run expected benefit. Tiers share source/link budgets, not value accounting.

\paragraph{Groups and deletion closures.}
Enumerate bounded, sealed, source-readable groups $X$ from the first gap,
within staging, permits and lifetime. Group gain can cross a zero-gain first
chunk. Deletion removes committed descendant closures, affecting all shared
consumers; candidate ancestors and read-pinned objects are excluded. In
tentative state $C'$, each closure's current density is
\begin{equation}
\rho(Z)=\frac{\Phi_r((C'\setminus Z)\cup X,t)-\Phi_r(C'\cup X,t)}{b(Z)},
\end{equation}
where $b$ counts distinct bytes. Greedily remove minimum-density closures
until the final state fits $\max(0,Q_r-S_r)$. Admit only if
\begin{equation}
\Phi_r(C_r,t)-\Phi_r((C_r\setminus Y)\cup X,t)>\eta\ge0.
\end{equation}
Choose minimum final harm, then fewer copied bytes. Bounded groups and greedy
closures can miss better combinations; no optimality or competitive ratio follows.

\begin{algorithm*}[tp]
\caption{\textsc{PrefixAdmission}: shared-prefix groups with staged replacement (revision 3)}
\label{alg:slack-landlord}
\textit{State:} $C$: committed object set; $Q$: granted bytes; $S$: staging reserve inside $Q$;
$\Phi$: current tier-local harm; $Y$: tentative deletion closures; $\eta\ge0$: improvement margin.
\par\medskip
\noindent\begin{minipage}[t]{.48\textwidth}\vspace{0pt}
\begin{tabular}[t]{@{}r@{\hspace{0.5em}}>{\raggedright\arraybackslash}p{\dimexpr\linewidth-2.5em\relax}@{}}
\multicolumn{2}{@{}l}{\textbf{Plan}$(i,r,t)$, independently for HBM and DRAM}\\
1 & Attach compatible committed shared prefix; $\mathcal F\gets\varnothing$.\\
2 & Pending group or unacknowledged reclaim $\Rightarrow$ \textsc{Pause}.\\
3 & \textbf{for} each bounded, sealed, source-readable group $X$ starting at the first missing object \textbf{do}\\
4 & $\quad$Require valid source lease, placement and path permits; finish estimate before expiry.\\
5 & $\quad$Require $b(X)\le S$ and $B_{physical}+b(X)\le Q$.\\
6 & $\quad C'\gets C$; $Y\gets\varnothing$.\\
7 & $\quad$\textbf{while} $b(C'\cup X)>\max(0,Q-S)$ \textbf{do}\\
8 & $\qquad Z\gets$ legal descendant closure minimizing current joint loss per distinct byte.\\
9 & $\qquad$Exclude candidate ancestors and read-pinned objects; no legal $Z\Rightarrow$ reject $X$.\\
10 & $\qquad C'\gets C'\setminus Z$; $Y\gets Y\cup Z$; recompute closure losses.\\
11 & $\quad$\textbf{end while}\\
12 & $\quad\Phi(C,t)-\Phi(C'\cup X,t)>\eta\Rightarrow\mathcal F\gets\mathcal F\cup\{(X,Y)\}$.\\
13 & \textbf{end for}; $\mathcal F=\varnothing\Rightarrow$ \textsc{Pause}.\\
14 & Choose minimum final harm, then minimum copied bytes; reserve $X$ within $Q$.\\
15 & Keep $C$ visible; copy $X$ from source as one PENDING group.\\
\end{tabular}\end{minipage}\hfill
\begin{minipage}[t]{.48\textwidth}\vspace{0pt}
\begin{tabular}[t]{@{}r@{\hspace{0.5em}}>{\raggedright\arraybackslash}p{\dimexpr\linewidth-2.5em\relax}@{}}
\multicolumn{2}{@{}l}{\textbf{Finish}$(X,Y,t')$}\\
16 & Validate all shard ACKs, value/route epoch, layout, deadline, quota and read leases.\\
17 & Recompute $\Phi(C,t')-\Phi((C\setminus Y)\cup X,t')>\eta$.\\
18 & \textbf{if} all checks pass \textbf{then} atomically publish $C\gets(C\setminus Y)\cup X$; enqueue $Y$ for reclaim.\\
19 & \textbf{else} retain old $C$; abort $X$ and enqueue fenced cancellation/reclaim.\\
20 & Retain physical charges and retry reclaim IDs until Worker free ACKs.\\
\multicolumn{2}{@{}l}{\textbf{OnEvent}$(e)$}\\
21 & New sealed KV, value/reference changes, free ACKs or restored permits trigger retry from the current gap.\\
22 & Request completion removes active value; optional zero-weight retention has a bounded TTL.\\
23 & Expired final references reclaim unpinned objects; shared live references remain.\\
24 & Hard foreground shrink aborts pending copies and revokes safe closures; physical release still requires ACK.\\
25 & Recovery reserves primary destinations before acquiring bounded source read leases.\\
26 & Lease expiry requests cancellation; only completion/cancel ACK permits physical free.\\
\end{tabular}\end{minipage}
\end{algorithm*}


Rejection pauses only protection. KV, value/reference, capacity or permit
changes retry from the first gap; missing source KV requires regeneration.
There is no permanent continuation ban. New rounds attach compatible resident
prefixes without copying; completing one user preserves other users' state.

\subsection{Commit, reclamation and recovery}
Algorithm~\ref{alg:slack-landlord} stages \emph{inside} $Q_r$ and revalidates
all-shard ACKs, epochs, layout, deadline, leases, quota and improvement before
atomically replacing the prefix. Failed copies retain old protection;
insufficient staging pauses admission. Physical accounting is
\begin{equation}
B_{physical}=B_{committed}+B_{copying}+B_{reclaiming}.
\end{equation}
Reclaim IDs retry until DMA/reads finish or cancel and physical free ACKs;
stale ACKs cannot free newer allocations. The reference permits one pending
group per pool and no further admission while reclaim is unacknowledged.

\begin{figure*}[tp]
\centering\includegraphics[width=\textwidth]{figures/commit-recovery.pdf}
\caption{Commit and recovery semantics common to both revisions. Failure
between tier ACKs leaves HBM chunks 0--1 and DRAM chunk 0. Recovery copies HBM
into primary buffers and recomputes suffix 2. C/X/dash mean committed/aborted/absent.
This schematic does not depict replacement or bound pressure races.}
\Description{A logical ACK timeline, distinct tier frontiers, and a physical
backup-to-primary copy followed by suffix recomputation.}
\label{fig:commit}
\end{figure*}

Recovery reserves primary destinations before source-read leases. Expiry
requests cancellation; it cannot free active DMA. Committed state is loaded
(Figure~\ref{fig:commit}); uncovered prompt/output positions are recomputed/replayed
before visible output resumes. Pending state is unrecoverable. Hard foreground
revocation may force fallback and is excluded from normal-replacement preservation.

\subsection{Algorithm 3: foreground-constrained HBM allowance}
\paragraph{Memory reservations.}
Before request/step/recovery reservations and watchdog events, account for
disjoint fixed memory $M_f$, primary reservations $M_a$, next incremental
allocation $M_n$, burst reserve $R_b$ and allocator margin $R_a$. For actual
allocator-lendable bytes $L_p$,
\begin{align}
H_p &= \pos{M_{physical}-M_f-M_a-M_n-R_b-R_a},\\
Q_{mem} &= \min(M_{registered},Q_{policy},L_p,H_p).
\end{align}
Unlendable arenas report $L_p=0$; CUDA free is insufficient. $R_b$ bounds growth
during stop/cancel/free ACK, using scheduler token/batch bounds; historical
quantiles give only empirical reserves. Display $\alpha=Q/M_{registered}$.

\paragraph{Interference constraint.}
Memory headroom does not isolate engines, bandwidth or locks. Each bounded
DMA batch checks expiring tokens for its source, target and links; both tiers
share bottleneck budgets. Batch size bounds non-preemptible service.

The default configurable target is $\epsilon=1\%$ overhead in no-failure p99
TTFT, p99 per-request TPOT and goodput loss versus matched backup-disabled
runs. One-sided confidence bounds must meet the target. Calibration needs
independent matched arrivals/tokens, cache and platform controls, correlation
and multiplicity handling, and held-out load regions. Certificates bind tested
joint quota/rate/batch domains and expiry, not arbitrary future bursts.

Unknown domains, stale samples, expired permits or foreground violations
revoke HBM allowance; free ACK still gates reuse. Production rejects modeled
certificates. Online monitoring triggers revocation without observing a
simultaneous no-backup counterfactual.

\paragraph{Enforcement.}
Quota is bounded by memory, certification and useful residency/copy demand.
Foreground events shrink it immediately; growth requires sustained valid
headroom, demand and completed reclaim, with a bytes-per-second limit.
Observation gaps reset stable time; epoch/sequence IDs make retries idempotent.
Actual foreground allocation checks physical feasibility, not published quota.
\begin{algorithm}[tp]
\caption{\textsc{ForegroundBudget}: revocable HBM allowance (revision 3)}
\label{alg:guarded-aimd}
\begin{tabular}{@{}r@{\hspace{0.5em}}>{\raggedright\arraybackslash}p{\dimexpr\columnwidth-2.5em\relax}@{}}
\multicolumn{2}{@{}l}{\textbf{BeforeReservation / Watchdog}$(p,o)$}\\
1 & Exact epoch/sequence retry returns the old decision; stale or conflicting IDs reject.\\
2 & Reconcile committed, copying and reclaiming bytes with physical inventory.\\
3 & Compute $Q_{mem}$ after disjoint current/next primary, burst and allocator reserves.\\
4 & Invalid samples/domain, expiry or foreground violation $\Rightarrow Q^*\gets0$; reset stable time.\\
5 & Otherwise $Q^*\gets\min(Q_{mem},Q_{cert},D)$; $D$ includes useful residency and staging.\\
6 & $Q^*<Q\Rightarrow Q\gets Q^*$ immediately; reset stable time.\\
7 & Valid sustained headroom, demand and no reclaim debt $\Rightarrow Q\gets\min(Q^*,Q+v_{grow}\Delta t)$.\\
8 & Publish $\alpha=Q/M_{registered}$ and expiring path/batch copy permits.\\
9 & Pause unsafe copies, abort/revoke safe objects, and retry reclaim IDs until free ACK.\\
10 & Foreground allocation requires actual physical feasibility, not a smaller published quota.\\
11 & Each bounded DMA batch rechecks permits and atomically debits all shared path budgets.\\
\end{tabular}
\end{algorithm}


Under truthful disjoint reservations, atomic allocator enforcement and a
valid bound on demand during cancellation, backup allocations cannot consume
reserved foreground capacity. This memory invariant is distinct from the
statistical interference target. Unbounded arrivals, missing cancellation
ACKs, or an allocator that cannot lend memory invalidate availability claims;
they do not authorize reusing an address before release.
